% Generated draft: requires mathematical and source review.
\documentclass[11pt]{article}
\usepackage[margin=25mm]{geometry}
\usepackage{fontspec}
\setmainfont{Latin Modern Roman}
\usepackage{amsmath,amssymb,mathtools,unicode-math}
\setmathfont{Latin Modern Math}
\usepackage{enumitem,needspace,xurl,hyperref}
\setlength{\parindent}{0pt}
\setlength{\parskip}{5pt}
\setlength{\emergencystretch}{3em}
\newcommand{\N}{\mathbb{N}}
\newcommand{\Z}{\mathbb{Z}}
\newcommand{\Q}{\mathbb{Q}}
\newcommand{\R}{\mathbb{R}}
\newcommand{\C}{\mathbb{C}}
\newcommand{\F}{\mathbb{F}}
\newcommand{\A}{\mathbb{A}}
\newcommand{\PP}{\mathbb{P}}
\newcommand{\E}{\mathbb{E}}
\newcommand{\eps}{\varepsilon}
\newcommand{\dd}{\,\mathrm{d}}
\begin{document}
\section*{UnsolvedMath ID: 30001662}
\textbf{Source problem code:} catalogue entry
\subsection{Definitions and mathematical statement}
For positive integers \(d\) and \(n\ge d+1\), let \(T\subset\R^d\) be a set of \(n\) distinct points. Interpret general position as requiring every subset of \(T\) with at most \(d+1\) elements to be affinely independent. For \(S\subseteq T\) with \(|S|=d+1\), its convex hull \(\operatorname{conv}(S)\) is a \(d\)-simplex. For any \(x\in\R^d\), define
\[
D_T(x)=\#\bigl\{S\subseteq T:|S|=d+1,\ x\in\operatorname{int}_{\R^d}(\operatorname{conv}(S))\bigr\}.
\]
Thus simplices are counted by their vertex subsets, and a simplex containing \(x\) only on its boundary is excluded. No general-position condition on \(T\cup\{x\}\) is imposed.

The extremal question is to determine, as a function of \(n\) and \(d\),
\[
M_d(n)=\max_{\substack{T\subset\R^d,\ |T|=n\\ T\text{ in general position}}}\ \max_{x\in\R^d}D_T(x).
\]
The parameter range above is the nontrivial positive-dimensional range; the source does not state parameter bounds or define general position explicitly. The displayed formulation interprets ``greatest possible'' as allowing both \(T\) and \(x\) to vary. If \(T\) is instead intended to be fixed, the corresponding target is \(\max_{x\in\R^d}D_T(x)\), which depends on \(T\). The supplied background also mentions an abstract version, but gives no definition from which to formulate it.
\subsection{Short English statement}
Determine the largest possible number of simplices with \(d+1\) vertices from an \(n\)-point configuration in general position in \(\R^d\) whose interiors contain one common point. Boundary containment does not count. The meanings of general position and the optimization over configurations require confirmation against the source.
\subsection{Sources}
\begin{itemize}
\item Local registry: \texttt{data/unsolvedmath/problems.json}, supplied record with permanent UnsolvedMath ID \(30001662\), title ``Maximum Simplicial Depth of Point Configurations,'' exact source code \texttt{catalogue entry}. No individual proposer is supplied; the recorded proposed year is 2011.
\item Supplied report citation: \emph{Combinatorics} (2011), Oberwolfach Reports, pp.~5--83, DOI 10.4171/owr/2011/01. \url{https://doi.org/10.4171/owr/2011/01}. The registry labels the statement \texttt{corrected\_verified} and reports that the official report was checked and the statement repaired to match its passage. This is supplied verification metadata, not an independent verification here.
\item Supplied literature reference: Leader--Tan, \url{https://arxiv.org/abs/0906.4027}.
\item Supplied literature reference: Matoušek--Wagner, \url{https://arxiv.org/abs/1102.3515}.
\item Literature assessment dated exactly \texttt{2026-08-21}: the supplied curation labels the problem partially solved and states that the geometric and abstract versions have strong asymptotic bounds and settled low-dimensional cases, while the exact higher-dimensional extremal value is not known. This assessment is supplied dated curation derived from the OWR material, not independently verified; its status and cited claims require independent checking.
\end{itemize}
\end{document}
